% 同一三点连接（1—2—3）的无向形式与两种定向：1->2->3 与 1->2<-3。
\begin{tikzpicture}[panelmark/.style={font=\figurefont\fontsize{8}{10}\selectfont,text=FigureInk,anchor=north,inner sep=0pt,outer sep=0pt},
>=Stealth,
  v/.style={circle,draw=FigureStroke,fill=FigureBlueFill,line width=.8pt,minimum size=6.5mm,inner sep=0pt},
  e/.style={draw=FigureStroke,line width=.8pt},
  a/.style={draw=FigureStroke,line width=.8pt,->,shorten >=1pt}]
\begin{scope}[local bounding box=panelrow]
\begin{scope}[local bounding box=panela]
  \node[v] (1) at (0,0) {$1$}; \node[v] (2) at (1.1,0) {$2$}; \node[v] (3) at (2.2,0) {$3$};
  \draw[e] (1)--(2)--(3);

\end{scope}
\begin{scope}[xshift=3.4cm,local bounding box=panelb]
  \node[v] (1) at (0,0) {$1$}; \node[v] (2) at (1.1,0) {$2$}; \node[v] (3) at (2.2,0) {$3$};
  \draw[a] (1)--(2); \draw[a] (2)--(3);

\end{scope}
\begin{scope}[xshift=6.8cm,local bounding box=panelc]
  \node[v] (1) at (0,0) {$1$}; \node[v] (2) at (1.1,0) {$2$}; \node[v] (3) at (2.2,0) {$3$};
  \draw[a] (1)--(2); \draw[a] (3)--(2);

\end{scope}
\end{scope}
% Subfigure marks: 8 pt, centred below each panel; common top 3 mm below this row.
\coordinate (panelrowbase) at ($(panelrow.south)+(0,-3mm)$);
\node[panelmark] at (panela.center |- panelrowbase) {(a)};
\node[panelmark] at (panelb.center |- panelrowbase) {(b)};
\node[panelmark] at (panelc.center |- panelrowbase) {(c)};
\end{tikzpicture}
